% =====================================================================
%  Detección de Rostros Sintéticos Generados por StyleGAN3
%  Entrega Final — Computación Paralela y Distribuida, Universidad LEAD
% =====================================================================
%
%  PARA COMPILAR CON LA PLANTILLA OFICIAL DE IEEE (recomendado en Overleaf):
%    1. Sustituir la línea \documentclass de abajo por:
%         \documentclass[conference]{IEEEtran}
%    2. Eliminar el bloque "CONFIGURACIÓN IEEE MANUAL" completo.
%    3. Sustituir el bloque de título por el formato \IEEEauthorblockN.
%
%  Tal como está, compila con pdflatex sin necesidad de IEEEtran.
% =====================================================================

\documentclass[10pt,twocolumn,letterpaper]{article}

\usepackage[utf8]{inputenc}
\usepackage[T1]{fontenc}
% El paquete de español de babel no está instalado en todos los entornos.
% Si lo tienes disponible (Overleaf lo trae), puedes sustituir esta línea por:
%   \usepackage[spanish,es-nodecimaldot]{babel}
\usepackage[english]{babel}
\usepackage{times}
\usepackage{graphicx}
\usepackage{amsmath}
\usepackage{amssymb}
\usepackage{booktabs}
\usepackage{multirow}
\usepackage{siunitx}
\usepackage{xcolor}
\usepackage{caption}
\usepackage[hidelinks]{hyperref}

% ---------- CONFIGURACIÓN IEEE MANUAL ----------
\usepackage[letterpaper,top=0.75in,bottom=1in,left=0.625in,right=0.625in,
            columnsep=0.25in]{geometry}
\usepackage{titlesec}
\titleformat{\section}{\normalfont\normalsize\scshape\centering}{\Roman{section}.}{0.5em}{}
\titleformat{\subsection}{\normalfont\normalsize\itshape}{\Alph{subsection}.}{0.5em}{}
\titleformat{\subsubsection}{\normalfont\small\itshape}{\arabic{subsubsection})}{0.5em}{}
\renewcommand{\thesection}{\Roman{section}}
\renewcommand{\thesubsection}{\Alph{subsection}}
\captionsetup[table]{labelsep=newline,textfont=sc,labelfont=sc,justification=centering}
\captionsetup[figure]{font=small,labelfont=small}
\setlength{\parindent}{1em}
\sisetup{output-decimal-marker={,},group-separator={\,},
         group-minimum-digits=4}

% ------------------------------------------------

\title{\vspace{-1.2em}\textbf{\large Detección de Rostros Sintéticos Generados por
StyleGAN3: Diagnóstico de Sesgos en Conjuntos Públicos y Diseño de un
Pipeline Acelerado en Cómputo Paralelo}\vspace{-0.3em}}

\author{
\begin{tabular}{c}
Jason Barrantes Sánchez \quad Melany Ramírez Anchía \quad
Walter Bowyer Carpenter \quad Mauro Espinoza Hernández \\[0.3em]
\small Bachillerato en Ingeniería en Ciencia de Datos, Universidad LEAD \\
\small San José, Costa Rica \\
\small \texttt{\{jason.barrantes, melany.ramirez, walter.bowyer,
mauro.espinoza\}@ulead.ac.cr}
\end{tabular}
}
\date{}

% Los nombres los fija babel al cargarse, así que se redefinen aquí,
% después del preámbulo, para que la sustitución tenga efecto.
\AtBeginDocument{%
  \renewcommand{\abstractname}{Resumen}%
  \renewcommand{\refname}{Referencias}%
  \renewcommand{\tablename}{Tabla}%
  \renewcommand{\figurename}{Figura}%
}

\begin{document}
\maketitle

% =====================================================================
\begin{abstract}
\small
Se presenta un sistema de detección de rostros sintéticos generados por
StyleGAN3, desarrollado sobre el clúster Kabré del Centro Nacional de Alta
Tecnología. El aporte principal del trabajo no es el clasificador en sí sino
el diagnóstico que lo precedió: un modelo Vision Transformer que alcanzaba
\SI{82}{\percent} de aciertos sobre el conjunto público de referencia
detectaba solo el \SI{23}{\percent} de los rostros que el mismo generador
producía en tiempo real. Mediante tres hipótesis falsables evaluadas
experimentalmente se descartó el sesgo de compresión ($d=0{,}093$) y el del
filtro de remuestreo (dispersión de \num{3} puntos porcentuales entre seis
filtros), y se confirmó que el conjunto público había sido generado con un
parámetro de truncamiento bajo, lo que produjo un atajo estadístico
explotable: la detección cae de \SI{100}{\percent} con $\psi=0{,}40$ a
\SI{23}{\percent} con $\psi=1{,}00$. Corregida la composición del conjunto y
la estrategia de entrenamiento, y sustituido el transformador de parches por
un conjunto de ConvNeXt-Tiny y Swin-T, el sistema alcanza
\SI{99{,}7}{\percent} de aciertos y AUC-ROC de \num{0.9999} sobre rostros
generados con semillas nunca vistas, empleando la quinta parte de los
parámetros. El pipeline de preprocesamiento paralelo sostiene
\num{2198}~im\'agenes/s con dieciséis procesos frente a
\num{511}~im\'agenes/s en ejecución secuencial.
\end{abstract}

\vspace{0.4em}
\noindent\textbf{\small Palabras clave---}{\small Vision Transformer, ConvNeXt,
detección de deepfakes, StyleGAN3, sesgo de conjunto de datos, cómputo
paralelo, Kabré.}
\vspace{0.8em}

% =====================================================================
\section{Introducción}

Los modelos generativos han alcanzado un grado de realismo que vuelve
inviable la verificación humana de autenticidad. StyleGAN3~\cite{karras2021}
introdujo una formulación equivariante a traslación y rotación que elimina
los artefactos de \emph{aliasing} característicos de sus predecesores,
produciendo rostros cuya distinción a simple vista resulta prácticamente
imposible.

Este trabajo aborda la clasificación binaria de rostros humanos: dada una
imagen RGB $I \in \mathbb{R}^{H \times W \times 3}$, determinar si
corresponde a una fotografía auténtica ($y=1$) o a una síntesis de
StyleGAN3 ($y=0$). En el marco del curso de Computación Paralela y
Distribuida, el problema se aborda además desde la perspectiva del
rendimiento computacional: se cuantifica el beneficio del paralelismo
multiproceso en el preprocesamiento y el de la aceleración por GPU en el
entrenamiento.

El hallazgo central del trabajo, sin embargo, fue metodológico. Un modelo
que parecía funcionar correctamente resultó estar explotando una
particularidad del conjunto de datos público en lugar de aprender las
características del generador. La sección~\ref{sec:diagnostico} documenta el
proceso de diagnóstico que reveló esa situación, que consideramos la
contribución más transferible de este trabajo.

\subsection{Contribuciones}

\begin{enumerate}\itemsep0.15em
\item La identificación y caracterización experimental de un sesgo de
truncamiento en un conjunto de datos público ampliamente disponible, con la
demostración de que un clasificador entrenado sobre él no generaliza al
generador que dice representar.
\item Un protocolo de evaluación que separa explícitamente la distribución de
entrenamiento de la de prueba mediante semillas disjuntas del espacio
latente, permitiendo medir generalización real.
\item La evidencia de que, para esta tarea, una arquitectura convolucional
supera a un transformador de parches con un tercio de los parámetros,
resultado contrario a la expectativa inicial del proyecto.
\item Un pipeline reproducible en infraestructura HPC, verificado mediante
regeneración completa desde cero.
\end{enumerate}

% =====================================================================
\section{Trabajo relacionado}

Wang \emph{et al.}~\cite{wang2020} demostraron que un clasificador
convolucional entrenado sobre imágenes de un único generador puede
generalizar a arquitecturas no vistas, sugiriendo la existencia de
artefactos comunes entre familias de modelos generativos. Nuestros
resultados matizan esa conclusión: la generalización falla cuando la
distribución de entrenamiento cubre solo una porción del espacio de salida
del generador, aun tratándose del mismo generador.

Los Vision Transformer~\cite{dosovitskiy2021} han mostrado desempeño
competitivo frente a las arquitecturas convolucionales en clasificación de
imágenes, y diversos trabajos han explorado su uso en detección de
contenido sintético argumentando que la autoatención captura relaciones
espaciales de largo alcance. La sección~\ref{sec:arquitecturas} presenta
evidencia en sentido contrario para esta tarea concreta.

En cuanto a infraestructura, el uso de Polars~\cite{polars} para
procesamiento columnar multihilo y de Dask~\cite{dask} para cómputo
distribuido perezoso está documentado en la literatura de ingeniería de
datos, aunque su aplicación a pipelines de visión orientados a detección de
contenido sintético permanece poco explorada.

% =====================================================================
\section{Metodología}

\subsection{Conjunto de datos}

Se emplea el conjunto público \emph{10000 Real vs Fake Faces (StyleGAN3)},
disponible en Kaggle bajo licencia CC-BY-NC-SA-4.0. Contiene
\num{20000} imágenes balanceadas (\num{10000} auténticas y \num{10000}
sintéticas) en formato PNG sin pérdida a resolución $256\times256$.

Conviene señalar una discrepancia detectada durante el trabajo: la
documentación del conjunto refiere una resolución nativa de
$1024\times1024$, pero la inspección de los archivos reveló que fueron
distribuidos ya reducidos a $256\times256$. Este detalle resultó relevante
para descartar el entrenamiento a mayor resolución
(sección~\ref{sec:resolucion}).

\subsection{Infraestructura de cómputo}

Todo el procesamiento se ejecutó en el clúster Kabré del CeNAT,
partición \texttt{nukwa-l40s}. La tabla~\ref{tab:infra} resume la
configuración empleada.

\begin{table}[t]
\caption{Configuración de la infraestructura}
\label{tab:infra}
\centering\small
\begin{tabular}{ll}
\toprule
\textbf{Recurso} & \textbf{Especificación} \\
\midrule
GPU & NVIDIA L40S, \SI{46}{\giga\byte} VRAM \\
CPU & Intel Xeon Silver 4416+, 20 núcleos \\
Núcleos asignados & 16 por trabajo \\
Memoria & \SI{64}{\giga\byte} \\
CUDA / cuDNN & 12.1 / 9.1 \\
Almacenamiento & Lustre, cuota de \SI{20}{\giga\byte} \\
Planificador & SLURM, límite de \SI{4}{\hour} por trabajo \\
\bottomrule
\end{tabular}
\end{table}

\subsection{Diseño del pipeline}

El pipeline se organiza en seis etapas. El inventario de metadatos se
construye con Polars, que aprovecha los núcleos disponibles sin
configuración explícita. El preprocesamiento --- decodificación,
redimensionado y normalización --- se reparte entre dieciséis procesos
mediante \texttt{ProcessPoolExecutor}, evitando la limitación del
\emph{Global Interpreter Lock} de CPython. Los tensores resultantes se
almacenan en fragmentos \texttt{.npy}.

Una decisión de diseño relevante fue almacenar las imágenes como enteros de
ocho bits sin normalizar, aplicando la normalización de ImageNet en tiempo
de carga. Frente al almacenamiento en punto flotante de 32 bits, esto
reduce el espacio en disco por un factor de cuatro
(\SI{12}{\giga\byte} a \SI{2{,}9}{\giga\byte}), lo cual resultó
determinante dada la cuota disponible. El costo computacional de normalizar
al vuelo es despreciable frente al de decodificar la imagen.

% =====================================================================
\section{Diagnóstico del sesgo de conjunto}
\label{sec:diagnostico}

\subsection{Observación inicial}

El primer modelo entrenado, un ViT-B/16 con ajuste fino completo, alcanzó
\SI{82}{\percent} de exhaustividad sobre las imágenes sintéticas del
conjunto de prueba. Al evaluarlo sobre \num{200} rostros generados en el
momento con el mismo modelo StyleGAN3-R y semillas nuevas, la detección
descendió a \SI{23}{\percent}: peor que decidir al azar.

Dado que ambos conjuntos provenían nominalmente del mismo generador, la
discrepancia solo podía explicarse por una diferencia sistemática entre las
imágenes del conjunto público y las generadas directamente. Se plantearon
tres hipótesis falsables.

\subsection{Hipótesis 1: sesgo de compresión}

Se comparó formato, modo de color, dimensiones, peso por píxel y tablas de
cuantización JPEG sobre \num{300} imágenes de cada clase. Todas resultaron
ser PNG sin pérdida en modo RGB a $256\times256$. La diferencia en peso por
píxel entre clases arrojó un tamaño de efecto de $d=0{,}093$, despreciable.
\textbf{Hipótesis descartada.}

\subsection{Hipótesis 2: sesgo de remuestreo}

Se observó que la energía en altas frecuencias de las imágenes generadas
($6{,}806$) quedaba fuera del rango de ambas clases del conjunto público
(auténticas $6{,}673$; sintéticas $6{,}606$), con un tamaño de efecto
$d\approx0{,}80$ respecto a estas últimas. La hipótesis era que el filtro
empleado al reducir de $1024$ a $256$ píxeles dejaba una firma distinta.

Se regeneraron las mismas \num{200} semillas aplicando seis filtros
distintos (Lanczos, bicúbico, bilineal, área, vecino más próximo y
reducción directa). Pese a que la energía en altas frecuencias varió entre
$6{,}381$ y $7{,}315$, la tasa de detección permaneció entre
\SI{18{,}5}{\percent} y \SI{21{,}5}{\percent}: una dispersión de tres puntos
porcentuales. \textbf{Hipótesis descartada.}

\subsection{Hipótesis 3: sesgo de truncamiento}

El parámetro $\psi$ de truncamiento controla cuánto se aleja una muestra del
centroide del espacio latente. Valores bajos producen rostros suavizados y
próximos al promedio; valores altos, rostros diversos y texturados.

Se generaron \num{200} rostros para cada uno de cinco valores de $\psi$,
evaluándolos con el mismo modelo. La tabla~\ref{tab:truncation} muestra el
resultado.

\begin{table}[t]
\caption{Detección según el truncamiento del generador}
\label{tab:truncation}
\centering\small
\begin{tabular}{ccc}
\toprule
$\boldsymbol{\psi}$ & \textbf{Detección} & \textbf{$P(\text{sintética})$ media} \\
\midrule
0,40 & \SI{100{,}0}{\percent} & 0,9996 \\
0,55 & \SI{97{,}5}{\percent}  & 0,9678 \\
0,70 & \SI{76{,}5}{\percent}  & 0,7624 \\
0,85 & \SI{49{,}5}{\percent}  & 0,4919 \\
1,00 & \SI{23{,}0}{\percent}  & 0,2436 \\
\bottomrule
\end{tabular}
\end{table}

La relación es monótona y el rango abarca \num{77} puntos porcentuales.
\textbf{Hipótesis confirmada.}

Para verificar que el generador era efectivamente el mismo, se compararon
los perfiles de potencia radial obtenidos por transformada de Fourier
bidimensional. La distancia $L_2$ normalizada entre las sintéticas del
conjunto público y las generadas resultó de \num{0.00733}, menor que la
distancia entre las dos clases del propio conjunto (\num{0.01055}). Mismo
generador, distinta configuración.

\subsection{Interpretación}

El conjunto público fue generado con truncamiento bajo, de modo que todas
sus muestras sintéticas comparten una suavidad característica. El
clasificador aprendió a reconocer esa suavidad --- una regla válida pero
incompleta --- en lugar de los artefactos del generador. Ante rostros
diversos, que se aproximan estadísticamente a fotografías reales, la regla
deja de aplicar.

Este tipo de sesgo es particularmente insidioso porque no es visible en
ninguna métrica de validación convencional: el modelo obtiene resultados
excelentes mientras se le evalúe sobre la misma distribución sesgada.

% =====================================================================
\section{Corrección y resultados}

\subsection{Reconstrucción del conjunto}

Se rediseñó el conjunto estratificando $\psi$ para concentrar la mitad de
las muestras sintéticas en el régimen difícil:
\SI{25}{\percent} en $[0{,}40, 0{,}70]$,
\SI{25}{\percent} en $[0{,}70, 0{,}85]$ y
\SI{50}{\percent} en $[0{,}85, 1{,}00]$, con media $\bar{\psi}=0{,}795$.

Se añadió augmentación aplicada por igual a ambas clases: espejo
horizontal, compresión JPEG de calidad variable entre 70 y 100, ajustes de
brillo, contraste y saturación, y desenfoque o realce. La compresión resultó
la componente más relevante: al almacenarse los fragmentos como PNG sin
pérdida, el modelo podía apoyarse en detalles de alta frecuencia que no
sobreviven a ninguna imagen del mundo real.

El protocolo de evaluación reserva dos conjuntos generados con semillas
disjuntas de las de entrenamiento (\num{2000000}--\num{2008999}):
\texttt{val\_hard} (semillas \num{2500000}+) para seleccionar el punto de
corte y detener el entrenamiento, y \texttt{test\_hard} (semillas
\num{2600000}+) que no participa en ninguna decisión. Ambos con
$\psi=1{,}00$ y balance exacto de \num{500} imágenes por clase.

\subsection{Evolución del sistema}

La tabla~\ref{tab:evolucion} resume las cinco iteraciones sobre el conjunto
de prueba difícil.

\begin{table}[t]
\caption{Evolución sobre \texttt{test\_hard} ($\psi=1{,}00$, semillas nuevas)}
\label{tab:evolucion}
\centering\small
\begin{tabular}{lccc}
\toprule
\textbf{Versión} & \textbf{Aciertos} & \textbf{AUC} & \textbf{Detecta sint.} \\
\midrule
v1 --- base            & ---   & ---     & \SI{23{,}0}{\percent} \\
v2 --- datos corregidos& ---   & ---     & \SI{60{,}9}{\percent} \\
v3 --- ajuste de LR    & \SI{89{,}8}{\percent} & 0,9534 & \SI{87{,}0}{\percent} \\
v4 --- augmentación    & \SI{92{,}0}{\percent} & 0,9679 & \SI{92{,}8}{\percent} \\
v5 --- tres ViT        & \SI{94{,}0}{\percent} & 0,9832 & \SI{96{,}2}{\percent} \\
\textbf{ConvNeXt+Swin} & \textbf{\SI{99{,}7}{\percent}} & \textbf{0,9999}
                       & \textbf{\SI{99{,}8}{\percent}} \\
\bottomrule
\end{tabular}
\end{table}

\subsection{Sobre la resolución de entrada}
\label{sec:resolucion}

Se consideró entrenar a $384$ píxeles bajo la hipótesis de que los
artefactos del generador residen en altas frecuencias. La idea fue
descartada por un argumento de coherencia: las imágenes del conjunto
público fueron distribuidas ya reducidas a $256\times256$, de modo que
ampliarlas no recupera detalle sino que lo interpola. Los rostros generados
por nosotros, en cambio, parten de $1024\times1024$ reales. A $384$ píxeles
las sintéticas tendrían detalle genuino frente al detalle inventado de las
auténticas, introduciendo exactamente el tipo de atajo que el trabajo buscaba
eliminar.

Se optó por $256$ píxeles, la resolución nativa del conjunto, evitando el
descarte del \SI{12}{\percent} de información que suponía reducir a $224$.
Al no ser $224$ la resolución de preentrenamiento de las codificaciones
posicionales del ViT, estas se interpolaron de una rejilla de $14\times14$ a
una de $16\times16$.

\subsection{Comparación de arquitecturas}
\label{sec:arquitecturas}

El resultado más inesperado del trabajo aparece en la
tabla~\ref{tab:arquitecturas}.

\begin{table}[t]
\caption{Desempeño individual por arquitectura sobre \texttt{test\_hard}}
\label{tab:arquitecturas}
\centering\small
\begin{tabular}{lccc}
\toprule
\textbf{Arquitectura} & \textbf{Parámetros} & \textbf{Aciertos} & \textbf{AUC} \\
\midrule
ViT-B/16        & \num{85{,}8}\,M & \SI{93{,}3}{\percent} & 0,9778 \\
Swin-T          & \num{27{,}5}\,M & \SI{97{,}3}{\percent} & 0,9995 \\
ConvNeXt-Tiny   & \num{27{,}8}\,M & \SI{99{,}6}{\percent} & 0,9999 \\
\bottomrule
\end{tabular}
\end{table}

ConvNeXt-Tiny supera al Vision Transformer por más de seis puntos
porcentuales empleando un tercio de los parámetros. La explicación
propuesta remite al mecanismo de cada arquitectura: el ViT particiona la
imagen en parches de $16\times16$ píxeles y proyecta cada uno a un vector,
promediando implícitamente la información interna del parche. Los rastros
que deja StyleGAN3 --- textura de piel, transición del cabello, patrones del
iris --- operan a una escala menor que ese parche, de modo que se atenúan
en la proyección. Las convoluciones jerárquicas, que recorren la imagen con
filtros pequeños y construyen representaciones progresivamente más
abstractas, preservan esa información.

\subsection{Conjunto de modelos y punto de corte}

Promediar las predicciones de ConvNeXt-Tiny y Swin-T, cada uno evaluado
sobre la imagen y su reflejo horizontal, eleva el AUC a \num{0.9999}.
Resulta notable que este conjunto de dos modelos
(\num{55}\,M de parámetros) supere al de cinco que incluía los
tres ViT (\num{313}\,M), que alcanzaba \SI{97{,}9}{\percent}.

El punto de corte se seleccionó sobre \texttt{val\_hard} tomando el centro
de la meseta de valores que quedan dentro del \SI{0{,}4}{\percent} del
óptimo, en lugar del máximo exacto. En iteraciones previas el máximo exacto
oscilaba entre $0{,}05$ y $0{,}94$ de una época a otra --- ruido, no
señal --- y llegó a producir un umbral que rendía peor que el valor por
defecto. El criterio de meseta arrojó $\tau=0{,}285$ sobre una región que
abarca el \SI{52}{\percent} de la rejilla explorada, indicando una
separación entre clases lo bastante nítida como para que la elección exacta
resulte poco crítica.

Aplicado a \texttt{test\_hard}, este umbral logra \SI{99{,}70}{\percent}
frente al \SI{99{,}90}{\percent} que habría obtenido el mejor umbral
posible seleccionado sobre el propio conjunto de prueba. Esa diferencia de
dos décimas es el costo metodológico de no inspeccionar el conjunto final,
y se reporta explícitamente.

% =====================================================================
\section{Análisis de rendimiento}

\subsection{Preprocesamiento paralelo}

Se midió el rendimiento del preprocesamiento sobre \num{2000} imágenes
variando el número de procesos, con tres repeticiones por configuración y
calentamiento previo de la caché del sistema de archivos. La
tabla~\ref{tab:throughput} recoge los resultados.

\begin{table}[t]
\caption{Rendimiento del preprocesamiento paralelo}
\label{tab:throughput}
\centering\small
\begin{tabular}{ccc}
\toprule
\textbf{Procesos} & \textbf{Tiempo (s)} & \textbf{Imágenes/s} \\
\midrule
1  & 3,913 & 511,1 \\
4  & 2,470 & 809,9 \\
8  & 1,354 & 1477,1 \\
12 & 1,061 & 1884,3 \\
16 & 0,910 & 2198,3 \\
\bottomrule
\end{tabular}
\end{table}

El rendimiento crece de forma monótona hasta un factor de $4{,}3$ con
dieciséis procesos. En la construcción del conjunto completo, la etapa de
preprocesamiento sostuvo \num{2396}~im\'agenes/s sobre
\num{9000} imágenes.

\subsubsection{Limitación de la medición}

Debe señalarse una limitación metodológica. La referencia secuencial
resultó inestable pese a tres estrategias correctivas sucesivas:
calentamiento completo de la caché, fijación a un solo hilo de las
bibliotecas numéricas mediante \texttt{OMP\_NUM\_THREADS} y variables
equivalentes, y toma del mínimo entre repeticiones. La medición con dos
procesos arrojó valores por debajo de la unidad, lo cual es
físicamente implausible en una carga sin dependencias entre elementos.

La causa más probable es que el trabajo se ejecuta bajo un \emph{cgroup} de
SLURM que restringe a dieciséis de los veinte núcleos físicos del nodo,
mientras el reparto interno entre el proceso padre y los hijos queda a
criterio del planificador del sistema operativo. En consecuencia, el
tiempo $T_1$ no corresponde a una ejecución estrictamente secuencial.

Por ese motivo se reporta el rendimiento absoluto, que es consistente entre
corridas, y se omite el cálculo formal de \emph{speedup} y eficiencia
paralela, que requeriría un $T_1$ fiable. Una medición rigurosa exigiría un
nodo en régimen exclusivo con afinidad de CPU fijada explícitamente.

\subsection{Cuello de botella}

El comportamiento observado sugiere que la carga está limitada por
entrada/salida más que por cómputo. Las imágenes residen en un sistema de
archivos Lustre compartido por red, y el rendimiento se aplana al
aproximarse a los doce procesos. La optimización natural consistiría en
precargar el conjunto a almacenamiento local del nodo antes de procesar.

\subsection{Entrenamiento en GPU}

El entrenamiento de ConvNeXt-Tiny sostuvo \num{614}~im\'agenes/s con
precisión mixta, completando quince épocas sobre \num{16800} imágenes en
\SI{7{,}9}{\minute}. Swin-T alcanzó \num{423}~im\'agenes/s y
diecisiete épocas en \SI{13{,}1}{\minute}. La generación de rostros con
StyleGAN3 sostuvo \num{36{,}8}~im\'agenes/s a resolución
$1024\times1024$.

% =====================================================================
\section{Reproducibilidad}

Los rostros sintéticos no se almacenan: se registran sus semillas y el
valor de $\psi$ empleado, lo que permite regenerarlos de forma determinista
sin coste de almacenamiento.

Esta propiedad se verificó de manera involuntaria pero concluyente. Una
operación de limpieza de disco eliminó los modelos entrenados y el conjunto
procesado. La reconstrucción completa --- redescarga del conjunto público,
regeneración de las \num{10000} imágenes sintéticas y reentrenamiento de
ambas arquitecturas --- tomó treinta y cinco minutos y devolvió resultados
idénticos hasta el cuarto decimal: AUC de \num{0.9999} y \num{0.9995} para
ConvNeXt y Swin respectivamente, umbral $\tau=0{,}285$, meseta de
\num{95} valores. Se documenta aquí porque constituye una verificación de
reproducibilidad más estricta que cualquier prueba planificada.

% =====================================================================
\section{Discusión y limitaciones}

El sistema alcanza \SI{99{,}7}{\percent} de aciertos, cifra que conviene
interpretar con precisión. Lo que el modelo distingue son \emph{rostros
generados por StyleGAN3-R frente a fotografías del conjunto FFHQ}. No es un
detector universal de contenido generado por inteligencia artificial: ante
salidas de modelos de difusión como Stable Diffusion o Midjourney, o incluso
de versiones anteriores de StyleGAN, cabe esperar un desempeño
sustancialmente menor.

Esa expectativa se sustenta en el propio hallazgo del trabajo. Si un cambio
en el parámetro de truncamiento del mismo generador bastó para reducir la
detección del \SI{82}{\percent} al \SI{23}{\percent}, un cambio de
arquitectura generativa producirá con seguridad un desplazamiento mayor. La
lección es general: los detectores de contenido sintético aprenden la firma
del generador con el que fueron entrenados, y la validación debe diseñarse
para exponer esa dependencia en lugar de ocultarla.

Como línea de trabajo futura, la más pertinente sería evaluar la
transferencia a otros generadores y, en su caso, entrenar sobre una mezcla
de arquitecturas generativas para inducir invarianza.

% =====================================================================
\section{Conclusiones}

Un clasificador con métricas aparentemente satisfactorias sobre un conjunto
público resultó estar explotando un artefacto de la construcción de ese
conjunto. El diagnóstico requirió formular hipótesis falsables y
descartarlas experimentalmente: dos resultaron incorrectas antes de
identificar la causa real.

La corrección del sesgo elevó la detección del \SI{23}{\percent} al
\SI{87}{\percent}, y el cambio de arquitectura la llevó al
\SI{99{,}8}{\percent}. Contra la expectativa inicial del proyecto ---
centrado en Vision Transformers desde su planteamiento --- la arquitectura
convolucional resultó marcadamente superior para esta tarea, con un tercio
de los parámetros.

El trabajo deja dos conclusiones transferibles. La primera es que las
métricas de validación no detectan sesgos compartidos entre los conjuntos de
entrenamiento y prueba; validar contra datos generados bajo condiciones
deliberadamente distintas es la única forma de exponerlos. La segunda es que
la elección de arquitectura debe responder a la naturaleza de la señal
buscada: cuando esa señal reside en texturas de escala fina, un
transformador de parches grandes la atenúa antes de poder analizarla.

% =====================================================================
\section*{Disponibilidad}

El código fuente, los cuadernos, los scripts de ejecución y los resultados
experimentales están disponibles en
\url{https://github.com/sklinderton/vit-real-vs-fake-faces}. Los conjuntos
procesados y los modelos entrenados se excluyen del repositorio por su
tamaño; el repositorio incluye los scripts y manifiestos de semillas que
permiten regenerarlos íntegramente.

% =====================================================================
\begin{thebibliography}{9}
\small

\bibitem{goodfellow2014}
I.~J. Goodfellow \emph{et al.}, ``Generative adversarial nets,'' in
\emph{Advances in Neural Information Processing Systems}, vol.~27, 2014.

\bibitem{karras2021}
T.~Karras, M.~Aittala, S.~Laine, E.~Härkönen, J.~Hellsten, J.~Lehtinen, and
T.~Aila, ``Alias-free generative adversarial networks,'' in \emph{Advances
in Neural Information Processing Systems}, vol.~34, 2021.

\bibitem{dosovitskiy2021}
A.~Dosovitskiy \emph{et al.}, ``An image is worth 16x16 words: Transformers
for image recognition at scale,'' in \emph{International Conference on
Learning Representations}, 2021.

\bibitem{liu2022}
Z.~Liu, H.~Mao, C.-Y. Wu, C.~Feichtenhofer, T.~Darrell, and S.~Xie, ``A
ConvNet for the 2020s,'' in \emph{Proceedings of the IEEE/CVF Conference on
Computer Vision and Pattern Recognition}, 2022, pp. 11976--11986.

\bibitem{liu2021swin}
Z.~Liu \emph{et al.}, ``Swin Transformer: Hierarchical vision transformer
using shifted windows,'' in \emph{Proceedings of the IEEE/CVF International
Conference on Computer Vision}, 2021, pp. 10012--10022.

\bibitem{wang2020}
S.-Y. Wang, O.~Wang, R.~Zhang, A.~Owens, and A.~A. Efros, ``CNN-generated
images are surprisingly easy to spot... for now,'' in \emph{Proceedings of
the IEEE/CVF Conference on Computer Vision and Pattern Recognition}, 2020,
pp. 8695--8704.

\bibitem{dask}
Dask Development Team, ``Dask: Library for dynamic task scheduling,'' 2016.
[En línea]. Disponible: \url{https://dask.org}

\bibitem{polars}
Polars Developers, ``Polars: Fast multi-threaded DataFrame library,'' 2023.
[En línea]. Disponible: \url{https://pola.rs}

\end{thebibliography}

\end{document}
